\documentclass[10pt,a4paper]{amsart}
\usepackage[margin=19mm]{geometry}
\usepackage{amsmath,amssymb,mathtools}
\usepackage[scr=rsfs,cal=euler]{mathalfa}
\usepackage{xfrac}
\usepackage[unicode,hidelinks,bookmarks=false]{hyperref}
\newcommand{\ie}{i.e.\ }
\newcommand{\cf}{cf.\ }
\newcommand{\resp}{respectively\ }
\newcommand{\Subsection}{\subsection}
\providecommand{\nobreakdash}{\nobreak\hbox{-}}
\setlength{\emergencystretch}{2em}
\multlinegap0pt
\usepackage{tikz}
\usetikzlibrary{cd}
\usepackage{mathtools}
\usepackage{amssymb}
\usepackage{tikz}
\newtheorem{prop}{Proposition}
\newtheorem{lem}{Lemma}
\title{Cohomology of the stratified fundamental group of curves}
\author{V.Q. Bao, P.H. Hai,, D.V. Thinh}
\date{Source: arXiv:2507.20020v2 (CC BY 4.0)}
\begin{document}
\maketitle

\noindent\emph{Source and licence.} Cohomology of the stratified fundamental group of curves. Version arXiv:2507.20020v2 (CC BY 4.0), distributed by arXiv under the Creative Commons Attribution 4.0 International licence, http://creativecommons.org/licenses/by/4.0/, as recorded in the arXiv licence metadata for this exact version and shown on the abstract page https://arxiv.org/abs/2507.20020v2. Full licence notice: \texttt{licenses/arxiv-2507.20020-CC-BY-4.0.txt}.\par\smallskip

\noindent\textit{Editorial scope.} Counted unit: the article's lem 5.13 (source0.tex lines 1672--1689). The article's complete original proof of it is delivered with it (source0.tex lines 1690--1738, one \texttt{proof} environment of 49 lines). No mathematics is written, repaired or re-derived: the statement and the proof are the article's own lines, in the article's own notation, and no retained line is substituted or re-flowed.  Retained uncounted as the source-local context the proof needs: lines 1605--1612; lines 1657--1660.  Nothing was omitted from any retained span, so the package carries no omission record.  No retained line contains a citation, so the wrapper carries no bibliography and no bibliography entry.  The retained lines contain the article's own diagram code, which the wrapper typesets with the standard tikz packages; no image file is delivered and the package declares no asset dependency.  Editorial formatting only, in addition: an AMS \texttt{amsart} preamble that restates the article's own declarations and macros used by the retained lines, namely the environments \texttt{prop}, \texttt{lem} on separate counters, together with the standard packages those lines need.  The retained environments print in this document as Proposition 1, Lemma 1 in order of appearance, whereas the article prints the same items with the numbering of its own full document.  The complete native source of the article as distributed in the arXiv e-print for this exact version is bundled unchanged as \texttt{sources/182256/source0.tex}.
\begin{prop}\label{exact sequence of En}
For any positive integers $m$ and $n$, there is a Frobenius-equivariant
exact sequence of vector bundles over $X$:
\[
0 \longrightarrow \mathcal E_n \longrightarrow \mathcal E_{m+n}
\longrightarrow \mathcal E_m \longrightarrow 0.
\]
\end{prop}

\begin{equation*}\tag{C}\label{condition_Cartier}
C_2(M_2)=0,\qquad C_3(M_3)\subset i_{1,3}(M_1),\qquad
\overline C_3:M_3/i_{2,3}(M_2)\xrightarrow{\sim}M_1.
\end{equation*}

\begin{lem}\label{lemma 5.13}\leavevmode
\begin{itemize}
\item[i)] For all $m,n\geq1$, there is an exact sequence
\[
0\longrightarrow M_n\xrightarrow{i_{n,n+m}}M_{n+m}
\xrightarrow{\mathrm{pr}_{n+m,m}}M_m\longrightarrow0.
\]
In particular, $\dim M_n=n$, and every $M_n$ contains a lift of $\omega$
under $\mathrm{pr}_{n,1}$ (with $\mathrm{pr}_{1,1}=\mathrm{id}$).
\item[ii)] Assume \eqref{condition_Cartier}. For $n>2$ and
$s\in M_n$ with $\mathrm{pr}_{n,1}(s)=\omega$, there exist
$s'\in M_{n-2}$ and $a\in k^\times$ such that
\[
\mathrm{pr}_{n-2,1}(s')=\omega,
\qquad C_n(s)=a\,i_{n-2,n}(s').
\]
\end{itemize}
\end{lem}

\begin{proof}
For i), tensor the sequence in Proposition \ref{exact sequence of En}
with $\Omega_X$ and take cohomology. The connecting morphism is
Cartier-equivariant. By Serre duality, the Cartier operator on
$\mathrm H^1(X,\Omega_X\otimes\mathcal E_n)$ is dual to Frobenius on
$\mathrm H^0(X,\mathcal E_n^\vee)\simeq k$, hence is bijective. Thus the
target of the connecting morphism has trivial nilpotent part. Since the
semisimple--nilpotent decomposition of a finite-dimensional vector space
with semilinear endomorphism is functorial for equivariant maps, passing to
nilpotent parts gives the asserted short exact sequence. Its dimension
statement follows from $M_1=k\omega$.

For ii), put $s_3=\mathrm{pr}_{n,3}(s)$. The projection
$\mathrm{pr}_{n,3}:\mathcal E_n\to\mathcal E_3$ is compatible with the
chosen Frobenius isomorphisms. Naturality of the Cartier map for this
bundle morphism therefore gives the commutative diagram
\[
\begin{tikzcd}
M_n \arrow[r,"C_n"] \arrow[d,"\mathrm{pr}_{n,3}"']
& M_n \arrow[d,"\mathrm{pr}_{n,3}"]\\
M_3 \arrow[r,"C_3"'] & M_3.
\end{tikzcd}
\]
Thus condition \eqref{condition_Cartier} gives, concretely,
\[
\mathrm{pr}_{n,3}\bigl(C_n(s)\bigr)
=C_3\bigl(\mathrm{pr}_{n,3}(s)\bigr)
=C_3(s_3)=a\,i_{1,3}(\omega),\qquad a\ne0.
\]
Since $\mathrm{pr}_{n,2}=\mathrm{pr}_{3,2}\mathrm{pr}_{n,3}$ and
$\mathrm{pr}_{3,2}i_{1,3}=0$, it follows that
$\mathrm{pr}_{n,2}(C_n(s))=0$. Applying i) with $m=2$ and with its
index $n$ replaced by $n-2$, we obtain
$C_n(s)=i_{n-2,n}(u)$ for some $u\in M_{n-2}$.
Finally, the restriction of $\mathrm{pr}_{n,3}$ to the subbundle
$\mathcal E_{n-2}\subset\mathcal E_n$ is the composite
\[
\mathcal E_{n-2}\xrightarrow{\mathrm{pr}_{n-2,1}}\mathcal E_1
\xrightarrow{i_{1,3}}\mathcal E_3.
\]
Consequently
\[
\mathrm{pr}_{n,3}i_{n-2,n}
=i_{1,3}\mathrm{pr}_{n-2,1}.
\]
Comparison with the preceding display gives
$\mathrm{pr}_{n-2,1}(u)=a\omega$.
Taking $s'=a^{-1}u$ proves ii).
\end{proof}

\end{document}
